\section{Introduction}

Large language models (LLMs) have rapidly evolved from conversational assistants into autonomous software engineering agents capable of planning, implementing, debugging, and deploying \xwz{increasingly}{} complex \xwz{software}{infrastructure} stacks \xwz{and infrastructure workflows}{}~\cite{jimenez2024swebench,liu2024agentbench,frontiereng2026}.
\xwz{Recent}{Modern} agentic \xwz{models have}{development frameworks} demonstrated \xwz{strong}{increasingly sophisticated} capabilities in repository-level code generation, multi-file reasoning, \xwz{iterative tool invocation, and workflow}{ tool} orchestration.
\xwz{However, despite these advances,}{Yet} the evaluation methodologies for \xwz{such}{these} systems remain \xwz{largely}{} rooted in a \emph{benchmark-driven paradigm} \xwz{consisting of}{:} isolated, static, and short-horizon tasks, where each evaluation \xwz{instance}{unit} is \xwz{executed independently}{ independent} and the agent \xwz{state}{} is reset after every attempt.

This \xwz{evaluation}{} paradigm \xwz{exhibits a fundamental mismatch with real-world production engineering environments}{ has a fundamental limitation: \emph{real production engineering is not a collection of independent tasks}}.
In practice, production \xwz{software and}{} infrastructure engineering involve long-horizon iterative \xwz{workflows spanning}{development,} multi-file and cross-module reasoning, dependency management, runtime debugging, \xwz{toolchain coordination}{tool orchestration}, deployment validation, continuous integration pipelines. \xwz{Real production systems must further operate under imperfect runtime conditions including}{} environment instability, resource constraints, \xwz{asynchronous execution, and cascading failure propagation}{ and recovery from cascading failures}.
\xwz{Existing static}{Static} benchmarks \xwz{fail to}{cannot} capture these dynamics---they \xwz{primarily}{} measure instantaneous coding \xwz{task completion rather than sustained engineering capability under evolving runtime conditions}{ability, not engineering evolution}. \xwz{Moreover, conventional benchmarks provide limited visibility into operational efficiency, execution stability, recovery behavior, and cost-performance tradeoffs, all of which are critical factors in practical deployment environments}{}.

As a result, strong benchmark performance does not necessarily imply \underline{practical engineering utility}.
Models \xwz{that perform well on}{optimized for} benchmark-style coding may still \warn{fail} in realistic production environments where runtime correctness, execution robustness, and adaptive recovery become critical.
\xwz{In many cases, failures originate not from insufficient code generation ability, but from unstable tool interaction, environment inconsistency, cascading execution errors, and inefficient runtime orchestration. Existing benchmark paradigms largely overlook these system-level behaviors because they evaluate isolated task completion rather than continuous execution under evolving runtime conditions.}{}
We \xwz{therefore }{} argue that the field \xwz{must}{needs to} move \emph{beyond static benchmarks} toward \xwz{production-grounded}{} \emph{\xwz{assessment}{evaluation} frameworks}---systems that \xwz{gauge}{assess} agentic models in realistic, multi-stage \xwz{engineering}{production} environments with \xwz{persistent runtime state,}{} controlled intermediate state injection, multi-dimensional metrics, and systematic failure diagnosis. \xwz{Such assessment should not only measure whether a task succeeds, but also characterize how the system behaves throughout the execution process}{}.

\xwz{To this end, we}{We} present \eva{} (Runtime Assessment of Models in Production), a production-grounded \xwz{assessment}{evaluation} framework built upon the \yatcc{} \xwz{integrated serving}{integrated} platform.
\eva{} provides a unified \xwz{runtime assessment infrastructure}{evaluation architecture} supporting heterogeneous LLM providers and agent SDKs, including OpenHands, Claude Code, Codex CLI, Kimi CLI, through standardized API access \xwz{and orchestration services}{} via \warn{\textsc{AIhub}}.
\xwz{Unlike conventional benchmark pipelines, \eva{} assesses models using}{It features} realistic long-horizon compiler construction workloads organized as serial dependency chains\xwz{, reflecting practical software engineering and infrastructure development processes}{}. \zzc{The framework further introduce \workload{}, a representative workload enabling systematic analysis of scaffolding effects and execution complexity under production-style workflows.}{\xwz{The framework further introduces}{with} two \xwz{workload}{} difficulty levels, \zzc{\Progressive{}}{\yatcc{}} and \zzc{\Incremental{}}{\yatcc{}-Hard}, \xwz{ enabling systematic analysis of scaffolding effects and execution complexity under production-style workflows}{a \resurrection{} mechanism that decomposes cascading failure into diagnostic components, and utility-oriented multi-dimensional metrics measuring both outcome quality and process efficiency}.} 
\xwz{A key design of RAMP is the Resurrection mechanism, which decomposes cascading workflow failures into recoverable diagnostic stages through controlled intermediate-state restoration. This mechanism substantially improves runtime observability and enables fine-grained analysis of downstream execution behavior that would otherwise be hidden by early-stage failures. In addition, RAMP introduces utility-oriented multi-dimensional metrics that jointly \warn{evaluate} outcome quality, execution efficiency, runtime cost, workflow robustness, recovery capability, and deployment-oriented operational behavior}{}.

% \textbf{Key Findings} from our production-oriented assessment include:
% \begin{itemize}
% \item Even the strongest model achieves only around 90\% task success rate, while the average success rate across models remains below 70\%.
% \item The cost-efficiency gap between models can exceed two orders of magnitude (up to 100×) under realistic production workloads.
% \item Up to 72\% of tasks that succeed on conventional benchmarks fail under production-scale workflow execution.
% \item Production execution introduces 3–5× higher failure rates than static benchmarks.
% \item More than 60\% of execution failures originate from system interaction rather than core reasoning.
% \item Long-horizon workflows reduce task completion rates by over 40\%.
% \item Workflow-level assessment reveals up to 2.8× performance divergence between models with similar benchmark scores.
% \item Runtime recovery capability improves production success rates by up to 35\%.
% \item Environment dependency issues account for over 45\% of software engineering task failures.
% \end{itemize}
%These findings collectively suggest that benchmark-centric evaluation substantially underestimates the complexity and operational challenges of real-world production environments.


In summary, the paper makes the following contributions:
\begin{enumerate}[nosep]
\item We propose \eva{} as a production-oriented \xwz{assessment}{evaluation} \emph{framework}, not merely another benchmark, that moves beyond synthetic isolated tasks toward persistent, runtime-observable evaluation of agentic models under realistic engineering environments.
\item We \xwz{construct}{design} realistic \xwz{compiler construction workloads}{long-horizon engineering workloads} with strict serial dependencies, verifiable intermediate states, and two difficulty levels that systematically quantify the impact of scaffolding and execution complexity in agentic software engineering tasks.
\item We introduce a unified \xwz{assessment infrastructure supporting heterogeneous LLMs and agent frameworks while introducing}{evaluation architecture and} utility-oriented metrics\xwz{ that characterize not only task outcomes,}{, measuring not just outcome scores} but \xwz{also runtime efficiency}{process cost}, failure patterns, and deployment-relevant \xwz{operational cost}{efficiency}.
\item We provide empirical evidence across 22 model configurations \zzc{and}{} \warn{four} agent backends, \zzc{}{and two workload difficulty level,} demonstrating that serial \xwz{assessment}{evaluation} reveals failure patterns invisible to static benchmarks and that resurrection-based recovery substantially improves downstream diagnostic observability.
\end{enumerate}